% Columnas en paralelo: dos órdenes cuyos términos se corresponden fila a fila.
% La correspondencia es `reciproca debil` (doble punta punteada), rotulada con su nombre.
\documentclass[tikz,border=12pt]{standalone}
\usepackage{diagrama}
\begin{document}
\begin{tikzpicture}[every node/.append style={font=\px{30}}]
  \node[titulo=44] at (4,5.4) {Título de la correlación};
  \node[font=\bfseries\itshape\px{36}] at (0.5,4.4) {Orden 1};
  \node[font=\bfseries\itshape\px{36}] at (7.5,4.4) {Orden 2};
  \foreach \i/\y/\a/\b in {1/3.4/Término 1/Término 1', 2/2.2/Término 2/Término 2', 3/1.0/Término 3/Término 3'} {
    \node (a\i) at (0.5,\y) {\a};
    \node (b\i) at (7.5,\y) {\b};
    \draw[reciproca debil] (a\i) -- node[rotulo] {\textit{Corresponde}} (b\i);   % e\i
  }
  \draw[derivacion] (a1) -- (a2);  % e4
  \draw[derivacion] (a2) -- (a3);  % e5
  \draw[derivacion] (b1) -- (b2);  % e6
  \draw[derivacion] (b2) -- (b3);  % e7
\end{tikzpicture}
\end{document}
